%! AP Statistics — Unit 1, Lesson 1.3 — Warm-Up
\documentclass[10pt]{article}
\usepackage{apstats-article}
\usepackage{apstats-boxes}

%=============================================================================
\begin{document}

\pageheader{Unit 1, Lesson 1.3}{Warm-Up}
\namedateperiod

\vspace{0.18in}

\noindent Below is a list of 20 students' favorite school subjects.

\vspace{0.1in}

\begin{center}
\small
\begin{tabular}{*{10}{c}}
Math & Science & English & Math & Art \\
Science & Math & History & English & Math \\
Art & Science & Math & English & Science \\
History & Math & Art & Science & English \\
\end{tabular}
\end{center}

\vspace{0.18in}

\begin{enumerate}

  \item Can you tell which subject is most popular by just scanning the list above?
        Why or why not?\\[6pt]
        \writeline\\[1pt]\writeline

  \item Use tally marks to count the number of students who chose each subject.
        Then write the count in the \textbf{Frequency} column.

\vspace{0.1in}
\renewcommand{\arraystretch}{1.8}
\begin{tabularx}{\linewidth}{@{} >{\bfseries}p{0.22\linewidth}
                                    C{0.26\linewidth} C{0.2\linewidth} @{}}
\rowcolor{sky}
\textbf{Subject} & \textbf{Tally} & \textbf{Frequency} \\
\midrule
Art     & & \\
English & & \\
History & & \\
Math    & & \\
Science & & \\
\midrule
\rowcolor{sky}\textbf{Total} & & \\
\end{tabularx}

  \item Does the total match the number of students in the list? \blank{1.5cm}
        Why is checking the total a good habit?\\[6pt]
        \writeline

  \item From Lesson 1.2: What \textit{type} of variable is ``favorite school subject''?
        \blank{3cm}\\[4pt]
        Is it meaningful to find the \textit{average} favorite subject? Explain.\\[6pt]
        \writeline

\end{enumerate}

\vspace{0.14in}

\begin{tcolorbox}[colback=greenbg, colframe=greenacc, boxrule=0.8pt, arc=2mm,
    left=4mm, right=4mm, top=3mm, bottom=3mm]
\small
\begin{itemize}[leftmargin=1.3em, itemsep=6pt]
  \item The table you just built is a \textbf{frequency table} --- it organizes
        categorical data into counts.
  \item Today we extend this to a \textbf{relative frequency table} by dividing each
        count by the total. This proportion is what allows fair comparison across groups.
  \item Describing \textit{how often each category occurs} is called describing the
        \textbf{distribution} of a categorical variable.
\end{itemize}
\end{tcolorbox}

\vspace{0.14in}

\noindent\textcolor{slate}{\small One question you have going into today's lesson, or something
you're not sure about yet:}\\[8pt]
\writeline

\end{document}
